\reporttitle{第一阶段执行日志与过程文档}


团队名称（待团队确认）　成员（待团队确认）　学校（待团队确认）\par
赛题为面向边缘AI大语言模型推理的RISC-V自定义vdot指令设计。本文初稿日期为2026年11月2日，阶段一原实验日期为2026年10月9日，最终工程版本按《最终版本确认》绑定。\par


\section{实验目的}


本阶段建立香山硬件仿真、NEMU参考模型与裸机程序之间的协同验证环境。以hello输出检查程序启动和串口链路，以向量加法检查向量状态、访存、执行、写回及提交过程，为后续自定义点积指令提供基线。\par


\section{环境与部署过程}


实验使用Windows下的WSL2 Ubuntu 22.04.5，处理器为Ryzen 7 7840H。原机器物理内存约27.7GiB，WSL配置20GB内存及16GB swap。硬件与仿真工作放在WSL原生文件系统。\par


{\small\setlength{\tabcolsep}{3pt}\begin{longtable}{>{\raggedright\arraybackslash}p{0.25\linewidth}>{\raggedright\arraybackslash}p{0.69\linewidth}}
\toprule
\textbf{组件} & \textbf{固定版本或设置}\\ \midrule\endhead
xs-env & \seqsplit{b0b506481e839a5df8a580e76f5aa246dc60fe7c}\\
原香山基线 & \seqsplit{16ae9ddcda54fc9a2fddffad73174cf793ac7814}\\
原NEMU基线 & \seqsplit{a48a11f208ff81ad765591d83271d1b339fb3861}\\
nexus-am & \seqsplit{086a238062fdfb3281b745a537b361983d0154b5}\\
JDK / Mill / firtool & 21.0.12 / 0.12.3 / 1.62.1\\
Verilator & 5.052，固定ea338be提交\\
RISC-V GCC / binutils & 11.4.0 / 2.38\\
DUT & MinimalConfig、单核、VLEN128\\
原波形构建 & EMU\_THREADS=1、EMU\_TRACE=fst\\
\bottomrule\end{longtable}}


部署依次完成WSL资源调整、固定仓库与递归依赖检出、工具版本核对、NEMU reference构建、香山仿真器构建、hello/vadd编译和差分运行。各步保存命令、标准输出、退出码、版本与SHA。环境配置详情与完整锁文件作为附件交付。\par


以下是既有脚本的操作顺序，实际操作应在已核对身份的专用输出目录执行。\par


\begin{Verbatim}[breaklines,breakanywhere,breaksymbolleft={},breaksymbolright={},fontsize=\small]
source team/scripts/activate.sh
bash team/scripts/check-environment.sh
bash team/scripts/stage1.sh nemu
bash team/scripts/stage1.sh emu
TEAM_ID=正式团队标识 bash team/scripts/stage1.sh workload
bash team/scripts/stage1.sh run
\end{Verbatim}


原始启动曾在PC 0x10000004发生寄存器差异。定位发现两端默认flash分别含3条和6条启动指令，随后由同一boot-flash.S生成12字节镜像，并显式传入共同flash、128MB RAM和固定reference。保存的失败与修复日志可说明启动条件对差分验证的影响。\par


\section{hello与vadd验证结果}


hello输出包含团队标识和WSL IPv4。历史记录使用riscv-vdot-team，正式录屏应呈现确认后的团队标识及相应真实输出。下文附真实历史输出，正式团队标识由成员确认后补入。\par


vadd程序设置SEW8、LMUL1、VL8，执行两次向量加载、一次vadd.vv和一次向量存储，逐字节对比标量参考。\par


{\small\setlength{\tabcolsep}{3pt}\begin{longtable}{>{\raggedright\arraybackslash}p{0.25\linewidth}>{\raggedright\arraybackslash}p{0.69\linewidth}}
\toprule
\textbf{用例} & \textbf{输出字节}\\ \midrule\endhead
正数 & 9 9 9 9 9 9 9 9\\
负数补码 & 254 252 250 248 246 244 242 240\\
边界回绕 & 128 127 0 0 0 0 254 254\\
\bottomrule\end{longtable}}


原运行同时具备HELLO\_VADD\_PASS cases=3 lanes=24、GOOD TRAP、启用DiffTest及exit0。程序整体统计为16557条指令、42962周期，包含启动、UART和校验。真实日志摘录见附录。\par


\section{波形与执行过程}


处理器完成取指和译码后，由重命名、调度和操作数准备驱动向量执行。两个vle8.v提供输入，vadd.vv完成按元素加法，结果经写回进入目的向量寄存器并在ROB提交。DiffTest在架构状态层与NEMU对应执行结果比较。\par


原波形提取的vadd提交PC为0x80000138，机器码为0x021101d7，vecwen=1、wdest=3。\par


{\small\setlength{\tabcolsep}{3pt}\begin{longtable}{>{\raggedright\arraybackslash}p{0.1\linewidth}>{\raggedright\arraybackslash}p{0.16\linewidth}>{\raggedright\arraybackslash}p{0.16\linewidth}>{\raggedright\arraybackslash}p{0.16\linewidth}>{\raggedright\arraybackslash}p{0.34\linewidth}}
\toprule
\textbf{用例} & \textbf{输入事件坐标} & \textbf{输出事件坐标} & \textbf{提交事件坐标} & \textbf{运算输出低64位}\\ \midrule\endhead
正数 & 19625 & 19626 & 19631 & \seqsplit{0x0909090909090909}\\
负数 & 25460 & 25461 & 25466 & \seqsplit{0xf0f2f4f6f8fafcfe}\\
边界 & 33009 & 33010 & 33015 & \seqsplit{0xfefe000000007f80}\\
\bottomrule\end{longtable}}


按原波形展开时，应展示时钟、valid/ready、提交PC和结果信号，并用游标标注表内事件。具体信号层级以原波形与事件提取脚本为准。表内三列事件数值为采样坐标，结果可说明运算单元输出早于架构提交的流水关系。\par




完整FST位于本机\seqsplit{/home/riscv/riscv-lab/stage1-evidence/hello-vadd.fst}，大小220295598字节。历史波形按2026年10月9日基线标注。最终0f环境的hello/vadd资格由R5组合验收另行绑定。\par


\section{结论及附件}


阶段一已完成部署、共同启动条件、3类向量加法用例与差分检查。提交报告应附操作日志索引、工具/子模块版本、源文件、反汇编、事件CSV及波形定位信息。本稿提供原波形提取表与事件时间线，团队标识待确认。\par


附件位置为\seqsplit{team/evidence/stage1-20261009}，包含\seqsplit{hello-vadd-summary.log}、工具版本、子模块锁、boot-flash反汇编、\seqsplit{vadd-commit-events.csv}及\seqsplit{vector-fix-alu-events.csv}。运行源码位于\seqsplit{team/workloads/hello-vadd}。\par


\href{https://github.com/cc77-builder/riscv/blob/eca41d7cd959f5f3c4f240e382919c492d1ccb96/team/docs/stage1.md}{阶段一原记录} · \href{https://github.com/cc77-builder/riscv/blob/eca41d7cd959f5f3c4f240e382919c492d1ccb96/team/docs/environment.md}{环境记录}\par
\section{实际执行日志摘录}
下列片段直接取自原始归档，日期及路径按原记录保留。完整日志随本文源文件的evidence目录提供。
\subsection{构建和运行时间记录}
\begin{itemize}
\item 2026年10月9日20:42:15，主动中断并恢复C++构建，将并行度调为4，已完成对象保留。
\item 2026年10月9日20:53:32，EMU构建结束，状态文件记录exit\_code=0。
\item 2026年10月9日20:55:26，workload与共同flash构建结束，记录\seqsplit{WORKLOAD\_BUILD\_PASS}。
\item 原正常运行至21:00:28结束，退出状态为0。
\end{itemize}
\subsection{程序构建}
\VerbatimInput[breaklines,breakanywhere,breaksymbolleft={},breaksymbolright={},fontsize=\small]{evidence/workload-build.log}
\subsection{差分运行与结果}
\VerbatimInput[breaklines,breakanywhere,breaksymbolleft={},breaksymbolright={},fontsize=\small]{evidence/hello-vadd-summary.log}
\section{从原波形提取的事件时间线}
\begin{center}
\begin{tikzpicture}[x=1cm,y=1cm,>=Stealth]
\draw[->](0,0)--(11,0) node[right]{相对周期};
\foreach \x/\t in {1/0,2.3/1,8.8/6}{\draw(\x,.15)--(\x,-.15)node[below]{\t};}
\draw[thick](1,0)--(1,1) node[above]{输入握手};
\draw[thick](2.3,0)--(2.3,1.7) node[above]{输出握手};
\draw[thick](8.8,0)--(8.8,1) node[above]{ROB提交};
\node[align=left] at (5,-1.2){正数用例的全局坐标依次为19625、19626、19631。\\结果低64位为\seqsplit{0x0909090909090909}。};
\end{tikzpicture}
\end{center}
图中坐标来自原FST提取CSV，为事件重绘示意。负数和边界用例的对应坐标列于正文表格，完整原波形以SHA绑定。
\section{日志附件索引}
evidence目录保留NEMU构建、EMU构建、workload构建、原正常运行、原启动失败及并行度调整日志，并包含工具版本、子模块锁、两份事件CSV及信号索引。原始日志用于追溯实际操作，本文摘录用于阅读。
